\answer{OP36}
记各段长度为 $D_1,\ldots,D_n$，其和为 $1$。正长度线段能够组成非退化多边形，当且仅当最大长度严格小于其余长度之和，因此本题条件是
\[
\max_iD_i<\frac12.
\]
必要性来自沿其余各边连接最大边两端时的严格三角不等式。充分性还可保证构型无自交。记最大边长为 $L$，总长为 $S$，假设 $2L<S$。对半径 $r\ge L/2$ 的圆，长度 $D_i$ 对应的较小圆心角为
\[
\theta_i(r)=2\arcsin\frac{D_i}{2r}.
\]
若在 $r=L/2$ 时这些角之和至少为 $2\pi$，随半径增大，角和连续降至零，故存在半径使角和恰为 $2\pi$。否则，将最大边对应的角改用优弧角 $2\pi-\theta_L(r)$：需要其余角之和等于 $\theta_L(r)$。两者之差在 $r=L/2$ 时为负，而充分大的 $r$ 下与 $(S-2L)/r$ 同号，因而为正。由连续性仍存在所需半径。两种情形都可依次在圆周上截取这些正角，再连接相邻顶点，得到边长合要求的凸多边形，因此无自交。

排序后的均匀切点联合密度为常数 $(n-1)!$。变换为段长后，密度仍在单纯形 $D_i>0$、$\sum D_i=1$ 上为常数，故各段长交换对称。第一段大于 $1/2$ 当且仅当所有切点都在 $1/2$ 右边，所以
\[
\Prob(D_i>1/2)=\Prob(D_1>1/2)=2^{-(n-1)}.
\]
总长为 $1$，不可能同时有两段大于 $1/2$；任一段恰等于 $1/2$ 的概率为零。于是
\[
\boxed{\Prob(\text{能够组成多边形})=1-\frac{n}{2^{n-1}}.}
\]
特别地，$n=3$ 时概率为 $1/4$，$n=4$ 时为 $1/2$。
